%! TEX root = ../master.tex
\begin{theorem}[Rotation of axes]\label{thm:rotation-of-axes}
  Rotating the axes through an angle $\theta$ gives new coordinates $(x', y')$
  related to the old ones by
  \[
    \begin{aligned}
      x &= x' \cos\theta - y' \sin\theta, &
      x' &= x \cos\theta + y \sin\theta, \\
      y &= x' \sin\theta + y' \cos\theta, &
      y' &= -x \sin\theta + y \cos\theta.
    \end{aligned}
  \]
  If $B \neq 0$, the rotation with
  \[
    \cot 2\theta = \frac{A - C}{B}, \qquad 0 < \theta < \frac{\pi}{2},
  \]
  removes the cross term from $Ax^2 + Bxy + Cy^2 + Dx + Ey + F = 0$. In
  particular $\theta = \pi/4$ when $A = C$.
\end{theorem}

\begin{proof}
  The two systems are inverse to each other, since the matrix of a rotation by
  $\theta$ has inverse the rotation by $-\theta$. Substitute the left-hand
  formulas into the quadratic part. Only $x^2$, $xy$ and $y^2$ produce an
  $x'y'$ term, and its coefficient is
  \[
    -2A \sin\theta\cos\theta + B(\cos^2\theta - \sin^2\theta)
    + 2C \sin\theta\cos\theta
    = B\cos 2\theta - (A - C)\sin 2\theta.
  \]
  This vanishes exactly when $\cot 2\theta = (A - C)/B$, and $\cot$ takes
  every real value once on $(0, \pi)$.
\end{proof}
